\chapter{The sieve of one denominator}\label{chap:sieve}

For each denominator~$b$ that the loop over the denominators visits and that passes the tests of
Chapter~\ref{chap:denominators}, the search looks for the points $(a : y : b)$ with
$\gcd(a, b) = 1$, $|a| \le H$ and $a/b$ in a search interval. ratpoints does this with a sieve. The
numerators become the bits of bit arrays of 64-bit words, packed by the stride of the 2-adic class
of~$b$. The rows of the moduli for~$b$, which say for which numerators the form can be a square
modulo the modulus (Chapter~\ref{chap:rows}), are and-ed into them, in a first stage over every bit
array and a second stage over the bit arrays that the first one leaves nonzero. Each numerator whose
bit survives is tested for common factors with~$b$, then against the primes of a third stage, and
then checked exactly; its points go to the callback (Chapter~\ref{chap:output}).

This chapter describes the model of this sieve, \lname{Ratpoints.Sieve.sift} and its parts, and its
contract (Definition~\ref{def:sift}): from a delivery whose flag of the stop is clear, the sieve
of~$b$ delivers, in increasing order, what the exact check finds for the numerators of the search
intervals that pass the 2-adic word of the class of~$b$ and survive the rows of all the moduli. The
block lemma (Theorem~\ref{thm:blocks}), and through it the main theorems, rests on this contract:
as every row and the 2-adic word let the numerator of every point through
(Chapters~\ref{chap:rows} and~\ref{chap:denominators}), the sieve delivers the points of the
brute-force reference over~$b$ with their numerators in the ranges.
Theorem~\ref{thm:sift_eq_checkDen} restates the contract, with the exact check on, as the check of a
plain sieve; no main theorem needs it. The moduli, the cut points of the stages and the tables of
rows come from Chapter~\ref{chap:moduli}, the search intervals and the ranges of their numerators
from Chapter~\ref{chap:region}.

\section{The sieve in the C code and in the model}\label{sec:sieveC}

In the C code the sieve of a denominator is \verb|sift| (\csource{find_points.c}{3442}{3617}), which
cuts the bits of each search interval into pieces for \verb|_ratpoints_sift0|
(\csource{sift.c}{498}{1305}), and each numerator that survives goes through \verb|accepted|
(\csource{sift.c}{438}{446}) to the exact check (\csource{find_points.c}{5250}{5338}). The model has
a component for each step, in the sections below, with the data and the loops of the default build.

A C \verb|long| is an integer or a natural number, with a lemma that it stays in range, except in
the test for common factors and the bits of an interval, which compute on 64-bit words. The
contracts use Lean core and \verb|Std| only; that of a component which finds points says that, from
a delivery whose flag of the stop is clear, it returns the reference of the delivery
(Definition~\ref{def:deliver}) of the list of the points it finds.

\begin{definition}[The counters of the sieve]\label{def:counters}
  \lean{Ratpoints.Sieve.Counters}
  \leanok
  The six counters that the sieve keeps for the correction of the moduli
  (Chapter~\ref{chap:moduli}), $64$-bit words that wrap (\csource{rp-private.h}{579}{581}): the
  words that the first stage sweeps, the bit arrays that survive it, the numerators extracted, those
  prime to the denominator, those that pass the third stage, and the calls of the sieve.
\end{definition}

\section{Bit arrays of words}\label{sec:bitArrays}

A bit array of the C code is a vector of words, with macros for its operations in each build
(\csource{rp-private.h}{113}{300}). In the model it is~$\mathrm{rbaPack}$ consecutive \verb|UInt64|
words of an array, one definition of each operation serving every build; the proofs read the bits
through \verb|BitVec.getLsbD|, whose lemmas Lean core has, and need no \verb|bv_decide|.

\begin{definition}[Bit arrays of words]\label{def:bitArrays}
  \lean{Ratpoints.rbaPack, Ratpoints.rbaLength, Ratpoints.Sieve.andBitArray, Ratpoints.Sieve.clearLowBits, Ratpoints.Sieve.clearHighBits, Ratpoints.Sieve.isZeroAt, Ratpoints.Sieve.getLsbD_clearLowBits, Ratpoints.Sieve.getLsbD_clearHighBits, Ratpoints.Sieve.ctz, Ratpoints.Sieve.ctz_spec}
  \leanok
  A bit array has $\mathrm{rbaPack} = 4$ words, $\mathrm{rbaLength} = 256$ bits, as in the default
  build (\csource{rp-private.h}{172}{172}, \csource{rp-private.h}{174}{175}); bit~$t$ of bit
  array~$j$ is bit $t \bmod 64$ of word $\mathrm{rbaPack} \cdot j + \lfloor t/64 \rfloor$. The
  operations: bit array~$j$ of an array of words and-ed with bit array~$x$ of a row, word by word
  (\verb|AND|, \csource{rp-private.h}{176}{176}); the lowest~$s$ bits of bit array~$0$, or the
  highest~$s$ bits of bit array~$j$, cleared for $s < \mathrm{rbaLength}$ (\verb|MASKL|,
  \verb|MASKU|, \csource{rp-private.h}{206}{213}, \csource{rp-private.h}{214}{221}), so that the low
  mask keeps bit~$t$ of word~$y$ exactly when $s \le 64 y + t$ and the high mask keeps it in a word
  $y < \mathrm{rbaPack}\,(j + 1)$ exactly when $64 y + t + s < \mathrm{rbaLength}\,(j + 1)$; and
  whether bit array~$j$ is zero (\verb|TEST|, \verb|TESTZ|, \csource{rp-private.h}{179}{202}). The
  lowest set bit of a word $x \ne 0$ comes from the loop of \verb|RP_CTZL| for a compiler without
  the builtin, which shifts~$x$ right while its lowest bit is clear
  (\csource{rp-private.h}{429}{444}); its contract: the count is the number of trailing zeros
  of~$x$.
\end{definition}

The proofs use of the width only $\mathrm{rbaPack} > 0$ and, for the packing, $\mathrm{rbaLength} =
2^{\mathrm{rbaShift}}$ with $\mathrm{rbaShift} < 32$, so they hold for every build
(Chapter~\ref{chap:ccode}).

\section{Reductions by a reciprocal}\label{sec:reciprocal}

ratpoints reduces modulo~$p$ by a multiplication with the reciprocal of~$p$ wherever that is exact:
the row indices of the first two stages, the residues of a denominator for its stride and for the
third stage, the offsets of the rows, the third stage, the first word of the forbidden divisors of
a block and the Jacobi test (Chapter~\ref{chap:denominators}). The reciprocal is a field of the
sieve entry (Chapter~\ref{chap:moduli}).

\begin{definition}[The reductions by a reciprocal]\label{def:reciprocal}
  \lean{Ratpoints.Sieve.magic, Ratpoints.Sieve.mulHi, Ratpoints.Sieve.mulMod, Ratpoints.Sieve.mulDiv, Ratpoints.Sieve.modMul}
  \leanok
  The reciprocal of~$p$ is $\mathrm{magic}(p) = \lfloor (2^{64} - 1)/p \rfloor + 1$. The high half
  $\mathrm{mulHi}(a, b)$ of the 128-bit product of two words, which the C code takes from gcc's
  \verb|__uint128_t| (\csource{rp-private.h}{404}{408}, \csource{rp-private.h}{409}{410}), comes in
  the model from the four products of their 32-bit halves, no partial sum leaving a word. The
  remainder by the reciprocal~$m$, \verb|RP_MULMOD|, is $\mathrm{mulHi}(m u \bmod 2^{64}, p)$, and
  the quotient, \verb|RP_MULDIV|, is $\mathrm{mulHi}(u, m)$, the arguments cast to words as the
  macros cast them; the builds that compile no 128-bit product (no such type, or
  \verb|RP_MULMOD_DIVIDE| or \verb|RP_NO_INT128| set) divide instead
  (\csource{rp-private.h}{413}{414}, \csource{rp-private.h}{415}{416}), with the same results
  wherever the multiplication is exact. The remainder of an integer~$a$, \verb|mod_mul|, is the
  remainder~$r$ of~$|a|$, and $p - r$ for negative~$a$ with $r \ne 0$ (\csource{sift.c}{243}{247}).
\end{definition}

\begin{lemma}[The reductions are exact]\label{lem:reciprocal}
  \uses{def:reciprocal}
  \lean{Ratpoints.Sieve.toNat_mulHi, Ratpoints.Sieve.mulMod_magic, Ratpoints.Sieve.mulDiv_magic, Ratpoints.Sieve.modMul_magic}
  \leanok
  For all words, $\mathrm{mulHi}(a, b) = \lfloor a b / 2^{64} \rfloor$. For $0 < p$ and
  $u p < 2^{64}$ the remainder with the reciprocal of~$p$ is $u \bmod p$; for $1 < p$ and
  $u p < 2^{64}$ the quotient is $\lfloor u/p \rfloor$; for $0 < p$ and $|a|\,p < 2^{64}$,
  \verb|mod_mul| gives the Euclidean remainder $a \bmod p$.
\end{lemma}

\begin{proof}
  \leanok
  With $a = 2^{32} a_1 + a_0$ and $b = 2^{32} b_1 + b_0$, the middle sum of the halves is below
  $3 \cdot 2^{32}$, and its high half is the carry. With $\mathrm{magic}(p)\, p = 2^{64} + \delta$,
  $0 \le \delta < p$, and $u = k p + s$, $s < p$, we have $\mathrm{magic}(p)\, u = k\, 2^{64} +
  (s\, 2^{64} + u \delta)/p$, the second part an integer below~$2^{64}$ as $u \delta < 2^{64}$: the
  high half of the product is~$k$, and its low half times~$p$ has the high half~$s$. The
  reciprocal~$2^{64}$ of~$1$ reads as~$0$ on a word, as in the C code, which spoils the quotient but
  not the remainder.
\end{proof}

The C code reduces so only values below~$2^{32}$ modulo moduli below~$2^{10}$, which the contracts
of the callers prove; it divides larger values with its helper \verb|mod|
(\lname{Ratpoints.Sieve.modLong}, the Euclidean remainder) or, in the third stage and its set-up,
with the operator~\verb|%|.

\section{The packing of the numerators}\label{sec:packing}

The 2-adic information gives for each class of the denominators modulo~$64$ the word~$w$ of the
numerators modulo~$64$ that a point can have (Chapter~\ref{chap:denominators}). In most classes the
set bits of~$w$ lie in one residue class $a_0 \bmod 2^k$, and ratpoints sieves only that class:
bit~$n$ of the sieve of~$b$ stands for the numerator $a_0 + 2^k n$, and bit~$t$ of the bit array
with the \emph{word number}~$i$ is the bit $n = \mathrm{rbaLength}\, i + t$. Each modulus~$p$ then
reads the row of $b\, 2^{-k} \bmod p$, at the bit array $(i + \mathrm{offset}) \bmod p$, with an
offset computed once per run for each packing and each modulus.

\begin{definition}[The packing]\label{def:packing}
  \uses{def:bitArrays, def:reciprocal, def:sieveEntry, def:rowStore, def:tableOfPrime}
  \lean{Ratpoints.Sieve.NumClass, Ratpoints.Sieve.numClasses, Ratpoints.Sieve.packPattern, Ratpoints.Sieve.classOffsets, Ratpoints.Sieve.classRows, Ratpoints.Sieve.fillBpList, Ratpoints.Sieve.numeratorsToBits}
  \leanok
  A class (\csource{rp-private.h}{494}{495}) has a pattern, a stride exponent~$k$, a
  residue~$a_0$ and the position of its row of offsets. The class of a word~$w$ of the 2-adic
  information (\csource{find_points.c}{1064}{1081}): for $w = 0$ the stride~$1$ and the pattern~$0$;
  otherwise, for the lowest set bit~$a$ of~$w$, the largest $k \le 6$ below $\mathrm{numStrides}$
  such that every set bit of~$w$ is $\equiv a \bmod 2^k$, $a_0 = a \bmod 2^k$, and the pattern whose
  bit~$t$ is bit $(a_0 + 2^k t) \bmod 64$ of~$w$ (\csource{find_points.c}{1082}{1094}); with the
  classes the word \verb|den_bits| of the classes whose word is not~$0$. The offsets
  (\csource{find_points.c}{1142}{1158}): for each entry of an odd modulus~$p$ the inverses
  $\mathrm{dinv}[k] = 2^{-k} \bmod p$ and $(2^k\,\mathrm{rbaLength})^{-1} \bmod p$, by halvings, and
  for each packing $(k, a_0)$ of a class with a nonzero pattern a row with the offset
  $a_0\,(2^k\,\mathrm{rbaLength})^{-1} \bmod p$ plus the bias of~$p$, the least multiple of~$p$ from
  $\mathrm{rowBias} = 2^{31}$ on (\csource{find_points.c}{1164}{1182}). The residues
  $b\,\mathrm{dinv}[k] \bmod p$ of a denominator for the moduli of the first two stages, in three
  branches by the size of~$b$ (\csource{find_points.c}{3639}{3679}). The bits
  $\lceil (\mathrm{low} - a_0)/2^k \rceil$ to $\lfloor (\mathrm{high} - a_0)/2^k \rfloor$ of an
  interval of numerators, by arithmetic shifts of signed 64-bit words, in a plain form and in a form
  that splits each end into its multiple of~$2^k$ and its rest, which cannot overflow
  (\csource{find_points.c}{3355}{3370}).
\end{definition}

\begin{lemma}[What the packing computes]\label{lem:packing}
  \uses{def:packing, lem:reciprocal, def:firstStage, def:rowStore, def:sieveEntry, def:tableOfPrime, lem:closedRow}
  \lean{Ratpoints.Sieve.NumClass.Packs, Ratpoints.Sieve.numClasses_spec, Ratpoints.Sieve.SieveEntry.PackedFor, Ratpoints.Sieve.classOffsets_spec, Ratpoints.Sieve.fillBpList_eq_map, Ratpoints.Sieve.toInt_numeratorsToBits, Ratpoints.Sieve.SieveEntry.admits_mul, Ratpoints.Sieve.SieveEntry.reads_closedRow}
  \leanok
  \begin{enumerate}
    \item The class of~$w$ \emph{packs}~$w$: it is \emph{well formed} ($k \le 6$,
      $k < \mathrm{numStrides}$, $a_0 < 2^k$), every set bit of~$w$ is $\equiv a_0 \bmod 2^k$, and
      bit~$t$ of its pattern is bit $(a_0 + 2^k t) \bmod 64$ of~$w$; bit~$j$ of \verb|den_bits| is
      set exactly when $j < 64$ and the word of~$j$ is not~$0$.
    \item For entries of odd moduli $1 < p \le \mathrm{maxPrimeEven}$ with their reciprocals, the
      inverse of~$\mathrm{rbaLength}$ and a bias in $p\ZZ \cap [\mathrm{rowBias}, \mathrm{rowBias}
      + p)$, well-formed classes and an array large enough, every class with a nonzero pattern gets
      a row whose offset~$o$ for each entry is \emph{packed for} it: $\mathrm{rowBias} \le o <
      \mathrm{rowBias} + 2p$, $o\, 2^k\,\mathrm{rbaLength} \equiv a_0 \bmod p$ and
      $\mathrm{dinv}[k]\, 2^k \equiv 1 \bmod p$ with $\mathrm{dinv}[k] < p$; nothing else changes.
    \item For entries with $0 < p \le \mathrm{maxPrimeEven}$, their reciprocals and
      $\mathrm{dinv}[k] < p$, the residues of a denominator are $b\,\mathrm{dinv}[k] \bmod p$.
    \item For a well-formed class the split form gives the bits $-\lfloor (a_0 - \mathrm{low})/2^k
      \rfloor$ and $\lfloor (\mathrm{high} - a_0)/2^k \rfloor$, and so does the plain form when
      $a_0 - \mathrm{low} < 2^{63}$ and $\mathrm{high} - a_0 \ge -2^{63}$.
    \item If the tables of inverses of the factors of an entry invert their units and
      $v u \equiv 1 \bmod p$, the row of~$u b$ admits~$u a$ exactly when the row of~$b$
      admits~$a$; so the closed form of the row of~$b'$ with $2^k b' \equiv b$
      (Chapter~\ref{chap:rows}), read at an offset~$o$ with $o\, 2^k\,\mathrm{rbaLength} \equiv
      a_0$, has at bit~$j$ the admission of $a_0 + 2^k j$ by the row of~$b$.
  \end{enumerate}
\end{lemma}

\begin{proof}
  \uses{lem:reciprocal}
  \leanok
  (1) The mask of the multiples of~$2^k$ is empty beyond $k = 6$, which ends the loop over the
  strides; the scatter puts bit~$a$ at $(a - a_0)/2^k$, and the doubling repeats the pattern.
  (2) A halving multiplies by $2^{-1}$ modulo~$p$; the rows occupy disjoint ranges of the array,
  one for each key $2^k + a_0$; the values reduced are below $64\, p$, so the reductions are exact
  (Lemma~\ref{lem:reciprocal}), and so in~(3). (4) The arithmetic shift by~$k$ is the floor of the
  quotient by~$2^k$, the and with $2^k - 1$ the Euclidean remainder. (5) The table of a prime is
  read at $a\, b^{-1}$, that of a prime power at $a\, b^{-1}$ or $b\, a^{-1}$, which~$u$ does not
  change. Read at the offset~$o$, bit~$j$ is the admission of $o\,\mathrm{rbaLength} + j$ by the
  row of~$b'$, and multiplied by the unit~$2^k$ this index becomes $a_0 + 2^k j$ modulo~$p$.
\end{proof}

So the packed sieve keeps exactly the numerators of the class $a_0 \bmod 2^k$ that a sieve without
the packing keeps, in the same order, and the pattern lets the bit~$n$ through exactly when~$w$ lets
$a_0 + 2^k n$ through (\lname{Ratpoints.Sieve.NumClass.Packs.classBit\_bits}); a numerator outside
the class is one that~$w$ rejects. The third stage is not packed.

\section{The first stage}\label{sec:firstStage}

The first stage and-s the rows of the first~$s_1$ moduli into every bit array of a piece, the bit
arrays with the word numbers $w_0, \dots, w_0 + n - 1$, from the pattern of the class. The C code
has two arms of it, which the preprocessor chooses by the number \verb|RATPOINTS_CHUNK| of
registers: for a chunk from~$2$ to~$16$ the chunked arm, which and-s the rows of all the moduli into
a chunk of bit arrays in registers; otherwise the arm that and-s one row at a time into the whole
piece. The model has both and chooses between them by the same test.

\begin{definition}[The first stage]\label{def:firstStage}
  \uses{def:bitArrays, def:reciprocal, lem:reciprocal, def:tableOfPrime}
  \lean{Ratpoints.chunk, Ratpoints.Sieve.SieveSpec, Ratpoints.Sieve.SieveSpec.Fits, Ratpoints.Sieve.andWords, Ratpoints.Sieve.andRows, Ratpoints.Sieve.chunkStarts, Ratpoints.Sieve.andRowsChunked, Ratpoints.Sieve.tailLegs, Ratpoints.Sieve.andRows_spec, Ratpoints.Sieve.andRowsChunked_spec}
  \leanok
  The chunk is~$16$, as in the default build (\csource{rp-private.h}{222}{224}). A sieve spec
  (\csource{rp-private.h}{378}{379}) holds a modulus~$p$, its row for the denominator at hand
  ($p + \mathrm{chunk} - 1$ bit arrays), the offset of the row, and a start and an end for the
  chunked arm. It \emph{fits} when $0 < p \le \mathrm{maxPrimeEven}$,
  $\mathrm{rowBias} \le \mathrm{offset} < \mathrm{rowBias} + 2p$, its end is~$p$, and bit array
  $x < p + \mathrm{chunk} - 1$ of its row is bit array $x \bmod p$. The test \verb|small| says
  whether the word numbers of the piece lie between $-\mathrm{rowBias}$ and $\mathrm{rowBias} -
  2\,\mathrm{maxPrimeEven}$; then the sums of a word number and an offset are reduced with the
  reciprocal, and otherwise by \verb|mod|. The arm one row at a time (\csource{sift.c}{943}{1022})
  sets every bit array to the class pattern and then, for each modulus, and-s bit array~$j$ with bit
  array $(w_0 + \mathrm{offset} + j) \bmod p$ of the row. The chunked arm sets the start of each
  row to $(w_0 + \mathrm{offset}) \bmod p$ (\csource{sift.c}{596}{606}); with no modulus in the
  stage ($s_1 = 0$) it writes the class pattern into every bit array, and otherwise it finds where
  the whole chunks end (\csource{sift.c}{608}{620}). A whole chunk is loaded into registers as the
  class pattern and-ed with the first row from its start on, and-ed with the other rows and stored,
  each start moving on by the chunk modulo~$p$; the $t < \mathrm{chunk}$ bit arrays after the whole
  chunks go in legs of~$8$, $4$, $2$ and~$1$ registers, one for each set bit of~$t$
  (\csource{sift.c}{915}{940}). The contract of each arm (of the chunked one under
  $\mathrm{chunk} \le 16$): for specs that fit, the reciprocals of their moduli and the bounds of
  \verb|small|, word $y < n\,\mathrm{rbaPack}$ of the piece is the class pattern and-ed, for each
  spec, with word $y \bmod \mathrm{rbaPack}$ of bit array $(w_0 + \mathrm{offset} + \lfloor
  y/\mathrm{rbaPack} \rfloor) \bmod p$ of its row.
\end{definition}

The chunked arm reads what the other arm reads: from the start $s = (w_0 + \mathrm{offset} + c)
\bmod p$ of the chunk at the bit array~$c$, bit array $s + j < p + \mathrm{chunk} - 1$ of the row
is bit array $(w_0 + \mathrm{offset} + c + j) \bmod p$. That the legs store each bit array of a
tail once is decided for each $t < 16$; the test that chooses the arm gives $\mathrm{chunk} \le
16$. The rows that the run reads fit (Chapter~\ref{chap:rows}). The masks of the ends of a piece
come after the first stage, as in the C code.

\section{The second stage and the extraction}\label{sec:extraction}

\begin{definition}[The second stage and the extraction]\label{def:extraction}
  \uses{def:bitArrays, def:reciprocal, def:firstStage, def:candidate, def:packing, def:deliver, lem:reciprocal, def:callback, def:checkNumWith, def:tableOfPrime, lem:checkNumWith, lem:deliverList, def:counters}
  \lean{Ratpoints.Sieve.secondStage, Ratpoints.Sieve.walkWord, Ratpoints.Sieve.siftBitArray, Ratpoints.Sieve.siftBitArrayWhole, Ratpoints.Sieve.scanPiece, Ratpoints.Sieve.secondStage_spec, Ratpoints.Sieve.walkWord_spec, Ratpoints.Sieve.siftBitArray_spec, Ratpoints.Sieve.scanPiece_spec}
  \leanok
  The second stage on a copy of the bit array with the word number~$i$: for each of its moduli in
  turn, while the copy is nonzero, the copy is and-ed with bit array $(i + \mathrm{offset}) \bmod p$
  of the row, the index reduced as in the first stage (\csource{sift.c}{1175}{1202}). The walk
  through a word~$x$ whose bit~$t$ stands for $a_0 + s t$: while $x \ne 0$, the numerator of the
  lowest set bit goes to the test of a candidate and, if it passes, to the exact check
  (Definition~\ref{def:candidate}), the walk returning if the callback has stopped the search;
  then~$x$ becomes $x \mathbin{\&} (x - 1)$ (\csource{sift.c}{262}{264},
  \csource{sift.c}{1227}{1255}). A bit array (\csource{sift.c}{1213}{1226},
  \csource{sift.c}{1257}{1264}): the second stage on the copy and, if the copy is still nonzero,
  the walk through each of its words, bit~$t$ of word~$\kappa$ standing for
  $a_0 + 2^k(\mathrm{rbaLength}\, i + 64\kappa + t)$. The scan of a piece
  (\csource{sift.c}{1061}{1082}): a sentinel of ones just past the piece; the zero bit arrays are
  stepped over, the scan stops at the sentinel, and each other bit array is sifted as above. The
  contracts, for specs that fit, the reciprocals of their moduli and the bounds of \verb|small|,
  from a delivery whose flag is clear: the copy after the second stage is the copy and-ed with the
  rows of all the moduli of the stage; the walk, and a bit array in either arm, deliver in
  increasing order the points of the numerators of the set bits (of~$x$, and of the words and-ed
  with the rows of the second stage) that pass the test of a candidate; and if, for a
  test~$\varphi$ of the bits, bit~$t$ of word~$y$ of the piece and-ed with the rows of the second
  stage is $t < 64 \wedge \varphi(\mathrm{rbaLength}\, w_0 + 64 y + t)$, the scan delivers in the
  same way the points of the numerators $a_0 + 2^k j$ of the bits~$j$ of the piece that
  pass~$\varphi$ and the test of a candidate.
\end{definition}

Stopping the second stage early loses nothing, since a zero bit array stays zero. A build with
\verb|USE_LONG_IN_PHASE_2| and-s one word at a time (\lname{Ratpoints.Sieve.secondStageWord}); the
model has both arms, chosen by the constant \lname{Ratpoints.useLongInPhase2}, false as in the
default build, and proves that they deliver the same points
(\lname{Ratpoints.Sieve.siftBitArray\_spec}). The walk ends since clearing the lowest set bit makes
the word smaller, the scan since the sentinel is not zero.

\section{The test of a candidate and the exact check}\label{sec:candidate}

\begin{definition}[The test of a candidate and the exact check]\label{def:candidate}
  \uses{def:bitArrays, def:reciprocal, def:deliver, def:callback, def:sieveEntry, def:tableOfPrime, def:counters}
  \lean{Ratpoints.Sieve.relprime, Ratpoints.Sieve.relprimeOdd, Ratpoints.Sieve.checkOf, Ratpoints.Sieve.fillChecks, Ratpoints.Sieve.stage3, Ratpoints.Sieve.accepted, Ratpoints.Sieve.checkPoint}
  \leanok
  The test for common factors on signed 64-bit words (\csource{sift.c}{318}{327}): for $m = 0$,
  whether $n = 1$; otherwise~$m$ replaced by~$-m$ when it is negative, false when~$m$ and~$n$ are
  both even, and else the binary algorithm on their odd parts: while $n \ne m$, with $d = m - n$ and
  the mask $d \gg 63$, $n$ becomes $n + (d \mathbin{\&} \mathrm{mask})$ and~$m$ becomes
  $(d \oplus \mathrm{mask}) - \mathrm{mask}$ shifted right by its trailing zeros; the result is
  whether~$m$ ends as~$1$ (\csource{sift.c}{328}{342}). What the third stage needs of the entry of a
  prime~$p$ for~$b$ (\csource{sift.c}{414}{428}): the entry of its table of inverses at $b \bmod p$,
  or~$0$ when $p \mid b$; the bias~$p H$ when $2 p H < 2^{32}$, else~$0$; the reciprocal and the
  table of~$p$; for the entries of the stage, a thunk made once per denominator
  (\csource{sift.c}{405}{413}). The third stage for~$a$: each prime whose inverse is not~$0$
  rejects~$a$ when its table says that~$f$ is not a square modulo~$p$ at the residue of
  $a\, b^{-1}$, reduced with the reciprocal after adding the bias if there is one, and by division
  otherwise (\csource{sift.c}{291}{301}). The test of a candidate: $a$ is prime to~$b$ and passes
  the third stage (\csource{sift.c}{438}{446}). The exact check
  (\csource{find_points.c}{5250}{5338}): unless~$\mathrm{noCheck}$, for the value~$v$ of Horner's
  rule at~$a$ on the products $c_k b^{d - k}$ (a thunk made once per denominator), times~$b$ for
  odd~$d$, so that $v = F_c(a, b)$ (Definition~\ref{def:binaryForm}), and
  $y = \lfloor \sqrt{v} \rfloor$: when $v = y^2$, the call with $(a : y : b)$ and then, unless that
  call stopped the search, $y = 0$ or~$\mathrm{noY}$, the call with $(a : -y : b)$;
  under~$\mathrm{noCheck}$ one call with $(a : 0 : b)$; the points mapped back to the caller's curve
  after a reversal.
\end{definition}

\begin{lemma}[The test of a candidate]\label{lem:candidate}
  \uses{def:candidate, def:callback, def:checkNumWith, def:sieveEntry, def:counters}
  \lean{Ratpoints.Sieve.relprime_spec, Ratpoints.Sieve.SieveEntry.IsPrimeEntry, Ratpoints.Sieve.all_admits_fillChecks, Ratpoints.Sieve.accepted_spec, Ratpoints.Sieve.acceptedPoints_eq_ite}
  \leanok
  For words $m > -2^{63}$ and $n > 0$ the test for common factors gives whether $\gcd(m, n) = 1$.
  A \emph{prime entry} is a sieve entry whose single factor is a table of the kind of a prime, of a
  modulus $1 < p \le 2^{32}$ ($p$ need not be prime), whose table of inverses holds a value in
  $[1, p - 1]$ at each nonzero residue below~$p$ (that these are inverses is not needed), and whose
  modulus and reciprocal are those of~$p$. For prime entries, $\gcd(a, b) = 1$ and $|a| \le H$, the
  third stage lets~$a$ through exactly when the rows of~$b$ in all the entries admit~$a$. For
  $0 < b < 2^{63}$ and $|a| < 2^{63}$ the test of a candidate is $\gcd(a, b) = 1$ followed by the
  third stage. So for prime entries in the third stage with their reciprocals, $0 < b < 2^{63}$ and
  $|a| \le H < 2^{63}$, the points that the extraction delivers for~$a$ are those of the check with
  the flags (\lname{Ratpoints.Sieve.checkNumWith}, the points of the exact check if
  $\gcd(a, b) = 1$) when the rows of the third stage admit~$a$, and none otherwise.
\end{lemma}

\begin{proof}
  \uses{lem:reciprocal, def:tableOfPrime}
  \leanok
  For positive words, $d = m - n$ is a word and the step gives $\min(m, n)$ and the odd part of
  $|m - n|$; the invariant keeps both values odd and positive, so that no step leaves the range of
  a \verb|long| and the arithmetic of the 64-bit words is that of the C code; their~$\gcd$ stays
  the same and their sum decreases, so that the loop ends at $m = n = \gcd$. For $m \ne 0$, $-m$ is
  a word, and removing the powers of~$2$ changes whether the~$\gcd$ is~$1$ only when both are even.
  In the third stage an entry whose modulus~$p$ does not divide~$b$ is not skipped, since its table
  of inverses is not~$0$ there; the biased value lies in $[H, 2pH - H]$, both reductions are exact
  (Lemma~\ref{lem:reciprocal}), and the residue read is where the row of~$b$ reads the table
  (Chapter~\ref{chap:moduli}). An entry with $p \mid b$ is skipped, and its row admits~$a$, since
  $p \nmid a$ by the~$\gcd$.
\end{proof}

The check with the flags is the check of the reference (\lname{Ratpoints.Sieve.checkNum}, what the
brute-force reference finds for one numerator, mapped back after a reversal), without the points
with $y < 0$ under~$\mathrm{noY}$, and the one survivor $(a, 0, b)$ under~$\mathrm{noCheck}$
(Chapter~\ref{chap:output}). The third stage of the run holds only primes of the table, also after a
correction of the moduli (Chapter~\ref{chap:moduli}), and these are prime entries.

\section{One denominator}\label{sec:siftDen}

\begin{definition}[The sieve of one denominator]\label{def:sift}
  \uses{def:firstStage, def:extraction, def:candidate, def:packing, def:bitArrays, def:deliver, def:bruteForce, lem:packing, lem:candidate, lem:reciprocal, lem:runRows, def:callback, def:checkNumWith, def:exactBounds, lem:deliverList, def:counters}
  \lean{Ratpoints.Sieve.siftPiece, Ratpoints.Sieve.siftBits, Ratpoints.Sieve.sift, Ratpoints.Sieve.siftPiece_spec, Ratpoints.Sieve.siftBits_spec, Ratpoints.Sieve.sift_spec, Ratpoints.Sieve.RowStore.Fits.fits_run_sift}
  \leanok
  One piece (\csource{sift.c}{498}{531}): the test \verb|small|, a fresh array with room for the
  sentinel, the first stage with the first~$s_1$ moduli in the arm that $1 < \mathrm{chunk} \le 16$
  chooses, the masks of the ends (\csource{sift.c}{1030}{1041}) and the scan with the moduli of the
  second stage. The bits from~$\mathrm{low}$ to~$\mathrm{high}$ of an interval
  (\csource{find_points.c}{3563}{3612}): the word numbers from
  $\lfloor \mathrm{low}/\mathrm{rbaLength} \rfloor$ to
  $\lfloor \mathrm{high}/\mathrm{rbaLength} \rfloor$ in pieces of at most~$A$ bit arrays, the first
  and the last masked. One denominator (\csource{find_points.c}{3442}{3488}): for each search
  interval the bits of its numerators within the height bound (Chapter~\ref{chap:region}); at the
  first interval that the height bound does not rule out (also one without numerators) the sieve
  specs (Definition~\ref{def:runTables}), the modulus, the offset of the row of the class and the
  row of the residue of~$b$ for the stride of each entry of the first two stages, the missing rows
  built into their tables (Chapter~\ref{chap:rows}); then the bits of the interval in pieces. Each
  returns when the callback stops the search. The contract of one denominator: for $A \ge 1$,
  $0 < b < 2^{63}$, $H < 2^{63}$, prime entries in the third stage, the reciprocals of all the
  entries, the residues $b\,\mathrm{dinv}[k] \bmod p$ of the entries of the first two stages, whose
  tables of inverses invert, tables of rows that fit them, and a delivery whose flag is clear: for
  every word~$w$ that the class packs, a row of offsets packed for the class, and the ranges of the
  numerators that the bounds of the search intervals give at~$b$, the delivery is that of the points
  of the check with the flags of the numerators $-H, \dots, H$ that lie in a range, pass~$w$ and
  survive the entries of all three stages, in increasing order; when the callback does not stop the
  search, the counters of the sieve grow by a closed form that depends neither on the callback and
  the flags nor on the tables of rows; and the tables of rows afterwards fit the entries again.
\end{definition}

The contracts compose from the bottom: the walk gives that of a bit array, the bit arrays that of
the scan, and the scan with either arm of the first stage and the masks that of a piece; the pieces
of an interval lie end to end. Lemma~\ref{lem:packing} turns the bits into numerators and the rows
for the stride into the rows of~$b$, and Lemma~\ref{lem:candidate} the test of a candidate into the
rows of the third stage. The counters (Definition~\ref{def:counters}), which feed the correction
of the moduli (Chapter~\ref{chap:moduli}), are counted before the exact check, as in the C code.

\section{The stages compute the plain sieve}\label{sec:siftTheorem}

The plain sieve, a sieve in its simplest form, models no C code.

\begin{definition}[The plain sieve]\label{def:plainSieve}
  \uses{def:bruteForce, def:sieveEntry}
  \lean{Ratpoints.Sieve.sieveDen, Ratpoints.Sieve.checkDen, Ratpoints.Sieve.sieveDen_spec, Ratpoints.Sieve.checkDen_spec}
  \leanok
  For a denominator~$b$, the height bound~$H$ and a list of sieve entries
  (Chapter~\ref{chap:moduli}): an array of $2H + 1$ Booleans, entry~$i$ for the numerator
  $a = i - H$, all true at the start, each and-ed in turn with the row of~$b$ in every entry at~$a$.
  Its contract: entry~$i$ is the survivor test \lname{Ratpoints.Sieve.survives} at $i - H$, the
  conjunction of the rows. The check of such an array takes the numerators whose entry is true, in
  increasing order, and tests and checks each as the brute-force reference does. Its contract: if
  entry~$i$ is $\varphi(i - H)$ for a test~$\varphi$ of the numerators, the result is what the
  brute-force reference (Definition~\ref{def:bruteForce}) finds over~$b$ and the numerators
  $-H, \dots, H$ that pass~$\varphi$.
\end{definition}

\begin{theorem}[The stages compute the plain sieve]\label{thm:sift_eq_checkDen}
  \uses{def:sift, def:plainSieve, def:deliver, lem:runRows, def:callback, def:exactBounds, def:reversal, def:counters}
  \lean{Ratpoints.Sieve.run_sift_eq_deliver_checkDen}
  \leanok
  Let $A \ge 1$, $0 < b < 2^{63}$ and $H < 2^{63}$; let the class pack a word~$w$ and have a row of
  offsets packed for it, the entries of the first two stages have tables of inverses that invert and
  tables of rows that fit them, the entries of the third stage be prime entries, and the bounds of
  the search intervals give the ranges of the numerators at~$b$. Then, with the exact check on, from
  a delivery whose flag is clear, \verb|sift| on~$b$, with the reciprocals of the entries and the
  residues of~$b$ for the stride, delivers the points that the check of the plain sieve with the
  entries of all three stages finds, restricted to the numerators that lie in a range and pass~$w$,
  mapped back after a reversal, and without those with $y < 0$ under~$\mathrm{noY}$.
\end{theorem}

\begin{proof}
  \uses{def:sift, def:plainSieve, def:checkNumWith, lem:checkNumWith}
  \leanok
  The check of the plain sieve gives what the brute-force reference finds for the numerators that
  survive all the entries (Definition~\ref{def:plainSieve}), and the sieve delivers the points of
  the check with the flags of the numerators in the ranges that pass~$w$ and survive all the
  entries (Definition~\ref{def:sift}). Without~$\mathrm{noCheck}$ this check is the check of the
  reference, mapped back and filtered by~$\mathrm{noY}$, and filtering the numerators filters the
  points of the reference by their numerators. A numerator survives a concatenation of lists of
  entries exactly when it survives each list (\lname{Ratpoints.Sieve.survives\_append}), so that
  the entries of the three stages together give the survivor test of all of them.
\end{proof}

\section{Where the model differs from the C code}\label{sec:sieveDeviations}

The \verb|Deviation:| paragraphs of the components record each difference, and
Chapter~\ref{chap:ccode} sorts them by kind. The high half of a 128-bit product comes from the
32-bit halves, and the lowest set bit from the loop for a compiler without the builtin. The test
for common factors and the bits of an interval compute on 64-bit words in two's complement with the
arithmetic shift, which the C standard leaves to the implementation and gcc defines so; the words
wrap where a \verb|long| would overflow only in the plain form of the bits, for an end
within~$a_0$ of~$-2^{63}$; there both the C code and the model use the split form. The comparison
$2pH < 2^{32}$ is made on natural numbers, for one of \verb|double|s, with the same outcome for
every $H < 2^{63}$. The flag of the stop is tested after each walk, bit array, piece and interval,
where the C code returns from within the walk.
